% Authoritative LaTeX quick start. Keep commands aligned with the English edition.
\chapter*{Quick start | 빠른 시작}
\addcontentsline{toc}{chapter}{Quick start | 빠른 시작}
\label{chap:quick_start}
\markboth{Quick start | 빠른 시작}{}
\begingroup
\setlength{\parskip}{3pt}
\textbf{목표: C 실행 파일로 한 방향을 계산하고, 12행짜리 Rayleigh LUT를 만든다.}
아래 예시는 에어로졸과 수중 입자를 사용하지 않아 대형 Mie 자료 다운로드 없이 시작할 수 있다. 운영체제에 맞는 준비 명령을 선택한 뒤 같은 터미널에서 이어서 실행한다. 이미 빌드했다면 실행 폴더로 이동해 2단계부터 시작한다. 자세한 설치는 \hyperref[chap:02_installation]{\ref*{chap:02_installation}장}을 참고한다.

\section*{1. 실행 준비}
\textbf{Linux / Ubuntu WSL:} GCC 11 이상, Make, Git이 필요하다. Ubuntu에서 도구가 없다면 먼저 \texttt{sudo apt-get update}와 \texttt{sudo apt-get install build-essential git}를 실행한다. 다음 빌드는 x86-64 AVX2 CPU용이다. 이전 x86-64 CPU는 \texttt{-march=x86-64-v3}를 \texttt{-march=x86-64}로 바꾼다. 이미 복제한 저장소에서는 \texttt{git clone}을 생략한다.

\begin{ManualCode}
git clone https://github.com/brtnt/OCRT.git
cd OCRT/code/OCRT_C
BUILD_FLAGS='-std=c11 -O3 -march=x86-64-v3 -ffp-contract=fast'
BUILD_FLAGS+=' -fassociative-math -fno-signed-zeros -fno-trapping-math'
BUILD_FLAGS+=' -DOCRT_FAST_KERNELS -fopenmp -Isrc'
make CFLAGS="$BUILD_FLAGS"
./build/ocrt --version
export OMP_NUM_THREADS=1
GAS_OFF=(--gas-column-h2o 0 --gas-column-o3 0 --gas-column-no2 0
         --gas-column-o2 0 --gas-column-co2 0 --gas-column-ch4 0)
mkdir -p ../../task/quick_start
\end{ManualCode}

\textbf{Windows PowerShell:} Git과 MSYS2 UCRT64 GCC가 설치되어 있다고 가정한다. GCC가 없으면 UCRT64 터미널에서 \texttt{pacman -S mingw-w64-ucrt-x86\_64-gcc}를 먼저 실행한다. 아래 PATH는 MSYS2 기본 설치 위치의 예시다. 이 매뉴얼의 실제 실행 확인 환경은 Linux/WSL이다.

\begin{ManualCode}
git clone https://github.com/brtnt/OCRT.git
Set-Location OCRT\code\OCRT_C
$env:Path = 'C:\msys64\ucrt64\bin;' + $env:Path
.\build_win.bat
.\build\ocrt_x86-64-v3.exe --version
$env:OMP_NUM_THREADS = '1'
$gasOff = @('--gas-column-h2o','0','--gas-column-o3','0',
  '--gas-column-no2','0','--gas-column-o2','0',
  '--gas-column-co2','0','--gas-column-ch4','0')
New-Item -ItemType Directory -Force ..\..\task\quick_start | Out-Null
\end{ManualCode}

\textbf{준비 완료 기준:} 빌드가 오류 없이 끝나고 \texttt{--version}이 출력된다. 버전 문자열은 오래된 문구를 유지할 수 있으므로 소스 버전 확인에는 커밋도 기록한다. 실행 폴더는 계속 \texttt{code/OCRT\_C}로 둔다. \texttt{inputs/}는 이 폴더를 기준으로 읽는다.

\textbf{읽기 안내:} OCRT는 \textbf{RAA=180도가 선글린트 가지}다. \hyperref[chap:01_concepts_geometry]{\ref*{chap:01_concepts_geometry}장}에서 위성 방위각 변환을 확인한다. 이 작은 격자는 동작 확인용이며 생산용 정확도를 검증한 격자가 아니다. 여러 파장\textperiodcentered 기압\textperiodcentered 풍속의 LUT와 해양 시뮬레이션은 \hyperref[chap:04_examples]{부록~\ref*{chap:04_examples}}, 결과의 단위와 열 정의는 \hyperref[chap:05_output_formats]{\ref*{chap:05_output_formats}장}으로 이어진다. 같은 파일명으로 다시 실행하면 예제 결과가 갱신되므로 보관할 결과는 새 폴더에 저장한다.

\clearpage
\section*{2. 한 방향 계산}
555 nm, SZA 40도, VZA 30도, RAA 90도, 기압 1013.25 hPa의 분자 산란을 계산한다. 기체 흡수는 끄며, 검은 경계이므로 해면 반사는 포함하지 않는다.

\textbf{Linux / WSL:}
\begin{ManualCode}
./build/ocrt --surface black --sza 40 --vza 30 --raa 90 \
  --wavelength 555 --pressure 1013.25 "${GAS_OFF[@]}" \
  > ../../task/quick_start/single.txt 2> ../../task/quick_start/single.log
echo "exit=$?"
head -n 5 ../../task/quick_start/single.txt
\end{ManualCode}

\textbf{Windows PowerShell:}
\begin{ManualCode}
& .\build\ocrt_x86-64-v3.exe --surface black --sza 40 --vza 30 --raa 90 `
  --wavelength 555 --pressure 1013.25 @gasOff `
  > ..\..\task\quick_start\single.txt 2> ..\..\task\quick_start\single.log
if ($LASTEXITCODE -ne 0) { throw 'OCRT failed; inspect single.log' }
Get-Content ..\..\task\quick_start\single.txt -TotalCount 5
\end{ManualCode}

\textbf{정상 실행 확인:} 종료 코드가 0이고 \texttt{single.txt}가 비어 있지 않아야 한다. 첫 세 값은 I/Q/U, 이후에는 \texttt{key=value} 결과가 나온다. 확인한 Linux/WSL 실행의 I/Q/U는 약 \texttt{0.03925161}, \texttt{0.00667091}, \texttt{-0.01252474}였으며 CPU와 빌드에 따라 마지막 자릿수는 달라질 수 있다. I/Q/U가 NaN/Inf이거나 종료 오류가 있으면 \texttt{single.log}를 확인한다. 이 대기 전용 예시에서 부가 수중 진단값은 NaN과 유효성 플래그 0으로 나올 수 있다. 표준에러에 진단 문구가 있다는 사실만으로 실패를 뜻하지 않는다.

\section*{3. 작은 Rayleigh LUT 생성}
거친 Fresnel 해면을 포함하고, 직달 선글린트 항을 제외한다. \textbf{단일 방향 옵션인 \texttt{--vza}, \texttt{--raa}를 이 명령에 추가하지 않는다.}

\textbf{Linux / WSL:}
\begin{ManualCode}
./build/ocrt --surface black_fresnel_ocean --wind-speed 5 \
  --sza 40 --wavelength 555 --pressure 1013.25 \
  --decouple-sunglint "${GAS_OFF[@]}" \
  --lut-vza-step 30 --lut-vza-max 60 --lut-raa-step 90 \
  --output-full-grid ../../task/quick_start/rayleigh.csv \
  2> ../../task/quick_start/rayleigh.log
echo "exit=$?"
wc -l ../../task/quick_start/rayleigh.csv
head -n 2 ../../task/quick_start/rayleigh.csv
\end{ManualCode}

\textbf{Windows PowerShell:}
\begin{ManualCode}
& .\build\ocrt_x86-64-v3.exe --surface black_fresnel_ocean --wind-speed 5 `
  --sza 40 --wavelength 555 --pressure 1013.25 --decouple-sunglint @gasOff `
  --lut-vza-step 30 --lut-vza-max 60 --lut-raa-step 90 `
  --output-full-grid ..\..\task\quick_start\rayleigh.csv `
  2> ..\..\task\quick_start\rayleigh.log
if ($LASTEXITCODE -ne 0) { throw 'OCRT failed; inspect rayleigh.log' }
(Import-Csv ..\..\task\quick_start\rayleigh.csv | Measure-Object).Count
Get-Content ..\..\task\quick_start\rayleigh.csv -TotalCount 2
\end{ManualCode}

\textbf{정상 실행 확인:} VZA 0/30/60도와 RAA 0/90/180/270도의 조합으로 \textbf{12개 데이터 행, 13개 열}이 생긴다. Bash의 \texttt{wc -l}은 헤더를 포함해 13, PowerShell의 데이터 행 수는 12가 정상이다. \texttt{rho\_I}, \texttt{rho\_Q}, \texttt{rho\_U}를 확인한다. Q/U는 음수일 수 있다.

\endgroup
